\documentclass[10pt,a4paper,onecolumn]{article}

\usepackage{amsmath,amssymb}
\usepackage{graphicx}
\usepackage[margin=2cm]{geometry}
\usepackage{hyperref}
\usepackage{cite}
\usepackage{enumitem}
\usepackage{booktabs}

\hypersetup{colorlinks=true,linkcolor=black,urlcolor=blue,citecolor=blue}
\setlength{\parskip}{0.3ex}

\title{\textbf{Early Detection of Deforestation in Protected Forests \\
Using U-Net and Sentinel-2 Imagery} \\
\small Idea Concept Paper --- Deforest.id}
\author{Research Team --- Deforest.id}
\date{}

\begin{document}
\maketitle
\thispagestyle{empty}

\paragraph{Abstract.}
We propose an automated deforestation alert system for Indonesian
protected forests (\textit{hutan lindung}) using a U-Net deep
learning model on 10~m Sentinel-2 imagery. The key innovation is
a fully autonomous pipeline: training labels are generated via
temporal NDVI differencing, eliminating manual annotation. An
integrated three-level early warning framework (Waspada/Siaga/Kritis)
with protected-zone alert amplification directly maps to national
and international reporting standards. Initial results demonstrate
0.73~IoU and 0.84~Dice on hold-out test scenes, with reliable
detection at 0.64~ha per chip. The system is designed for rapid
deployment across Indonesia's $\sim$94~Mha forest estate without
ground surveys, supporting the FOLU Net Sink 2030 target \cite{menlhk2022folu}.

\section{Problem Statement}
Indonesia has the third-largest tropical forest area globally
($\sim94$~Mha), yet loses $\sim300$~kha of natural forest annually
\cite{hansen2013high,fao2022sofo}. Protected forests (\textit{hutan lindung}) --
which regulate watersheds, prevent erosion, and maintain soil
fertility for millions of people -- are especially vulnerable to
small-scale deforestation that evades coarse-resolution monitoring
systems. The cumulative impact is severe: patches smaller than
2~ha contribute over 56\% of tropical forest carbon emissions
\cite{xu2026nature}, while fragmentation exponentially accelerates
edge effects and biodiversity loss \cite{taubert2018fragmentation}.

\paragraph{The gap.}
Current operational systems cannot address this scale mismatch.
Global Forest Watch operates at 30~m resolution (Landsat), missing
sub-0.9~ha clearings. Indonesia's SIMONTANA system uses a minimum
mapping unit of 6.25~ha for REDD+ reporting
\cite{gfw2024simontana,unfccc2022indonesia} -- three times the 2~ha threshold above
which ecological impacts become severe. Crucially, no existing
system provides \textit{automated alert prioritization} for
protected zones, which legally require stricter oversight under
Indonesian law (UU~No.~41/1999). Norway's NICFI program \cite{nicfi2020} provides
high resolution imagery but only at monthly intervals and without
automated analysis. The core problem is a monitoring resolution
gap: deforestation occurring at 0.5--2~ha scales is simultaneously
too small for coarse sensors and too large to ignore ecologically.

\section{Research Description}
We present an end-to-end pipeline that combines Sentinel-2
multi-spectral imagery (10~m, 5-day revisit, open access) \cite{drusch2012sentinel2} with a
U-Net deep learning architecture \cite{ronneberger2015unet}. The pipeline is fully autonomous:
it generates its own training data, detects clouds without relying
on defective QA60 bands, and outputs area-based severity alerts
directly mappable to policy thresholds.

\subsection{Autonomous Training via Pseudo-Labeling}
The core technical innovation is a temporal NDVI differencing
method that generates high-quality training labels without any
manual annotation. For each location, we compute:
\begin{equation}
\Delta\text{NDVI}(x,y) = \text{NDVI}_{\text{pre}} - \text{NDVI}_{\text{post}}
\end{equation}
and threshold to produce binary deforestation masks:
\begin{equation}
\text{mask}(x,y) = 
\begin{cases}
1 & \Delta\text{NDVI} < -0.15 \\
0 & \text{otherwise}
\end{cases}
\end{equation}
Morphological closing ($3\times3$ kernel) removes salt-and-pepper
noise. This approach exploits the fact that deforestation causes
a sharp, sustained drop in NDVI that is distinct from seasonal
variation (which tends to be gradual and bidirectional). The
threshold $-0.15$ was calibrated against visual interpretation of
50 random samples across 3 scenes. This pseudo-labeling enables
scalability: a model trained on one region can be fine-tuned for
another without costly field campaigns.

\subsection{Early Warning Framework}
Each Sentinel-2 pixel represents $10\times10$~m = 0.01~ha, enabling
direct area conversion: $A(\text{ha}) = N_{\text{pixels}} \times 0.01$.
Detected deforested areas are classified into three severity levels
(Table~\ref{tab:thresholds}), with alerts automatically amplified
by one level for protected forest zones.

\begin{table}[ht]
\centering
\caption{Alert thresholds synthesizing FAO, ecological, and national
standards.}
\label{tab:thresholds}
\small
\begin{tabular}{lcc}
\toprule
Level & Area (ha) & Pixels \\
\midrule
\textit{Waspada} & $\ge 0.64$ & $\ge 64$ \\
\textit{Siaga}   & $\ge 2.00$ & $\ge 200$ \\
\textit{Kritis}  & $\ge 6.25$ & $\ge 625$ \\
\bottomrule
\end{tabular}
\end{table}

The 0.64~ha baseline aligns with FAO's forest definition minimum
(0.5~ha), 2~ha with the ecological fragmentation threshold
\cite{xu2026nature}, and 6.25~ha with Indonesia's REDD+ minimum
\cite{unfccc2022indonesia}. For \textit{hutan lindung}, area-based
amplification applies: \textit{Siaga} escalates to \textit{Kritis},
\textit{Waspada} escalates to \textit{Siaga}. This ensures protected
areas receive faster enforcement response commensurate with their
legal protection status.

\subsection{Cloud Detection Without QA60}
Raw Sentinel-2 QA60 bands from GEE contain only placeholder zeros.
We implement a simple but effective heuristic:
\[
\text{cloud}(x,y) = R{>}180 \land G{>}180 \land B{>}180
\land \text{NDVI}{<}0.10
\]
This captures thick cumulus clouds; residual thin clouds and
shadows are handled by the Dice loss's robustness to label noise
during training.

\subsection{Implementation Pipeline}
The pipeline operates in six stages:
\begin{enumerate}[nosep,itemsep=1pt]
    \item \textbf{Acquisition}: 7 Sentinel-2 scenes (L1C) from
          Riau and Central Kalimantan (2019--2024), downloaded
          via Google Earth Engine.
    \item \textbf{Tiling}: Scenes split into $512\times512$ blocks,
          then $64\times64$ chips with 50\% overlap, producing
          40,484 baseline chips and 20,103 deforest chips.
    \item \textbf{Cloud masking}: Heuristic filter applied per
          chip; cloud-contaminated chips flagged for exclusion.
    \item \textbf{Pseudo-labeling}: $\Delta\text{NDVI}<-0.15$ with
          morphological closing ($3\times3$) generates masks.
    \item \textbf{Training}: U-Net (4 encoder/decoder stages,
          32--512 filters, 5-channel input), Dice loss, Adam
          ($1\times10^{-4}$), batch 64, 50 epochs, split
          11,500/2,908/5,695. NIR band normalized to $[0,1]$
          via division by 10,000.
    \item \textbf{Inference}: Sliding window (50\% overlap),
          connected-component labeling, area aggregation, alert
          classification per threshold.
\end{enumerate}

\section{Initial Results}
Training on a single RTX~3060 (12~GB) required $\sim$2.7~min/epoch;
50 epochs completed in ${<}$3~h with num\_workers=0 for WSL
compatibility. On 5,695 hold-out chips from 2 unseen scenes:
\begin{itemize}[nosep]
    \item \textbf{IoU:} 0.7256
    \item \textbf{Dice:} 0.8410
    \item \textbf{Precision:} 0.8340
    \item \textbf{Recall:} 0.8480
\end{itemize}
The system reliably detects deforestation patches at 64~pixels
(0.64~ha), operating at our minimum alert threshold. By comparison,
a simple NDVI temporal threshold ($\Delta\text{NDVI}<-0.15$) alone
achieves only $\sim$0.53~F1 on the same test set, confirming that
the U-Net learns spatial context beyond spectral change.

\section{Impact and Scalability}
The proposed system would provide Indonesia's Ministry of
Environment and Forestry with automated, daily deforestation alerts
at 0.01~ha effective resolution -- a 900-fold improvement over the
6.25~ha REDD+ minimum. Three features ensure practical viability:

\paragraph{Zero annotation cost.}
The pseudo-labeling pipeline generates training labels from any
Sentinel-2 image pair, enabling scaling to all 38 Indonesian
provinces without field surveys. Estimated label generation cost
per province: $\sim$30 min of cloud compute for image pair
processing.

\paragraph{Commodity hardware feasibility.}
Training completes in under 3 hours on a single consumer GPU
(RTX~3060, $\sim$\$250). Inference on a full Sentinel-2 scene takes
${<}$5~minutes. This eliminates dependency on cloud compute or
specialized infrastructure, making the system accessible to
provincial forestry offices.

\paragraph{Policy alignment.}
Alert thresholds directly map to national REDD+ reporting (6.25~ha),
FAO forest definition (0.5~ha), and the 2~ha ecological threshold.
Protected-zone amplification requires no legislative change --
it operationalizes existing legal protection under UU~No.~41/1999.
The system directly supports the FOLU Net Sink 2030 target by
enabling earlier enforcement interventions.

\section{Roadmap}
\begin{enumerate}[nosep,itemsep=1pt]
    \item \textbf{Sentinel-1 SAR integration}: C-band radar to
          penetrate year-round cloud cover in Sumatra and
          Kalimantan, improving revisit from 5-day to near-daily.
    \item \textbf{Field validation}: Ground-truth campaigns at 3
          protected forest sites (Riau, Central Kalimantan, East
          Kalimantan) in collaboration with local forestry agencies.
    \item \textbf{Web platform}: Real-time dashboard with GEE
          streaming, on-the-fly inference, and email/WhatsApp alert
          notifications for district forestry officers.
    \item \textbf{Policy integration}: Direct linkage with KHLK
          (Kawasan Hutan Lindung dan Konservasi) boundary maps
          enabling automated severity amplification and quarterly
          reporting to KLHK.
\end{enumerate}

\small
\bibliographystyle{ieeetr}
\bibliography{icp-refs}
\end{document}
